\section{AMC 12 12A 2005}
        \begin{enumerate}
        	\item (\textit{AMC 12 12A 2005}) Two is $10 \%$ of $x$ and $20 \%$ of $y$. What is $x - y$?


 $\textbf{(A) } 1 \qquad \textbf{(B) } 2 \qquad \textbf{(C) } 5 \qquad \textbf{(D) } 10 \qquad \textbf{(E) } 20$


	\item (\textit{AMC 12 12A 2005}) The equations $2x + 7 = 3$ and $bx - 10 = - 2$ have the same solution for $x$. What is the value of $b$?


 $\textbf{(A) } -8 \qquad \textbf{(B) } -4 \qquad \textbf{(C) } -2 \qquad \textbf{(D) } 4 \qquad \textbf{(E) } 8$


	\item (\textit{AMC 12 12A 2005}) A rectangle with a diagonal of length $x$ is twice as long as it is wide. What is the area of the rectangle?


 $\textbf{(A) } \frac{1}{4}x^2 \qquad \textbf{(B) } \frac{2}{5}x^2 \qquad \textbf{(C) } \frac{1}{2}x^2 \qquad \textbf{(D) } x^2 \qquad \textbf{(E) } \frac{3}{2}x^2$


	\item (\textit{AMC 12 12A 2005}) A store normally sells windows at $\$100$ each. This week the store is offering one free window for each purchase of four. Dave needs seven windows and Doug needs eight windows. How many dollars will they save if they purchase the windows together rather than separately?


 $\textbf{(A) } 100 \qquad \textbf{(B) } 200 \qquad \textbf{(C) } 300 \qquad \textbf{(D) } 400 \qquad \textbf{(E) } 500$


	\item (\textit{AMC 12 12A 2005}) The average (mean) of $20$ numbers is $30$, and the average of $30$ other numbers is $20$. What is the average of all $50$ numbers?


 $\textbf{(A) } 23 \qquad \textbf{(B) } 24 \qquad \textbf{(C) } 25 \qquad \textbf{(D) } 26 \qquad \textbf{(E) } 27$


	\item (\textit{AMC 12 12A 2005}) Josh and Mike live $13$ miles apart. Yesterday Josh started to ride his bicycle toward Mike's house. A little later Mike started to ride his bicycle toward Josh's house. When they met, Josh had ridden for twice the length of time as Mike and at four-fifths of Mike's rate. How many miles had Mike ridden when they met?


 $\textbf{(A) } 4 \qquad \textbf{(B) } 5 \qquad \textbf{(C) } 6 \qquad \textbf{(D) } 7 \qquad \textbf{(E) } 8$


	\item (\textit{AMC 12 12A 2005}) Square $EFGH$ is inside the square $ABCD$ so that each side of $EFGH$ can be extended to pass through a vertex of $ABCD$. Square $ABCD$ has side length $\sqrt{50}$, $E$ is between $B$ and $H$, and $BE = 1$. What is the area of the inner square $EFGH$?

                    \begin{figure}[H]
                        \centering
                        \includegraphics[width=0.5\textwidth]{Problem/AMC_12_12A_2005_7_Problem_Figure_1.png}
                    \end{figure}
                    
$\textbf{(A) } 25 \qquad \textbf{(B) } 32 \qquad \textbf{(C) } 36 \qquad \textbf{(D) } 40 \qquad \textbf{(E) } 42$


	\item (\textit{AMC 12 12A 2005}) Let $A$, $M$, and $C$ be digits with


 
 \[(100A+10M+C)(A+M+C) = 2005.\]


 What is $A$?


 $\textbf{(A) } 1 \qquad \textbf{(B) } 2 \qquad \textbf{(C) } 3 \qquad \textbf{(D) } 4 \qquad \textbf{(E) } 5$


	\item (\textit{AMC 12 12A 2005}) There are two values of $a$ for which the equation $4x^2 + ax + 8x + 9 = 0$ has only one solution for $x$. What is the sum of those values of $a$?


 $\textbf{(A) } -16 \qquad \textbf{(B) } -8 \qquad \textbf{(C) } 0 \qquad \textbf{(D) } 8 \qquad \textbf{(E) } 20$


	\item (\textit{AMC 12 12A 2005}) A wooden cube $n$ units on a side is painted red on all six faces and then cut into $n^3$ unit cubes. Exactly one-fourth of the total number of faces of the unit cubes are red. What is $n$? 


 $\textbf{(A) } 3 \qquad \textbf{(B) } 4 \qquad \textbf{(C) } 5 \qquad \textbf{(D) } 6 \qquad \textbf{(E) } 7$


	\item (\textit{AMC 12 12A 2005}) How many three-digit numbers satisfy the property that the middle digit is the average of the first and the last digits?


 $\textbf{(A) } 41 \qquad \textbf{(B) } 42 \qquad \textbf{(C) } 43 \qquad \textbf{(D) } 44 \qquad \textbf{(E) } 45$


	\item (\textit{AMC 12 12A 2005}) A line passes through $A(1,1)$ and $B(100,1000)$. How many other points with integer coordinates are on the line and strictly between $A$ and $B$?


 $\textbf{(A) } 0 \qquad \textbf{(B) } 2 \qquad \textbf{(C) } 3 \qquad \textbf{(D) } 8 \qquad \textbf{(E) } 9$


	\item (\textit{AMC 12 12A 2005}) In the five-sided star shown, the letters $A$, $B$, $C$, $D$ and $E$ are replaced by the numbers $3$, $5$, $6$, $7$ and $9$, although not necessarily in that order. The sums of the numbers at the ends of the line segments $\overline{AB}$, $\overline{BC}$, $\overline{CD}$, $\overline{DE}$ and $\overline{EA}$ form an arithmetic sequence, although not necessarily in that order. What is the middle term of the arithmetic sequence?


 
                    \begin{figure}[H]
                        \centering
                        \includegraphics[width=0.5\textwidth]{Problem/AMC_12_12A_2005_13_Problem_Figure_1.png}
                    \end{figure}
                    
$\textbf{(A) } 9 \qquad \textbf{(B) } 10 \qquad \textbf{(C) } 11 \qquad \textbf{(D) } 12 \qquad \textbf{(E) } 13$


	\item (\textit{AMC 12 12A 2005}) On a standard die one of the dots is removed at random with each dot equally likely to be chosen. The die is then rolled. What is the probability that the top face has an odd number of dots?


 $\textbf{(A) } \frac{5}{11} \qquad \textbf{(B) } \frac{10}{21} \qquad \textbf{(C) } \frac{1}{2} \qquad \textbf{(D) } \frac{11}{21} \qquad \textbf{(E) } \frac{6}{11}$


	\item (\textit{AMC 12 12A 2005}) Let $\overline{AB}$ be a diameter of a circle and $C$ be a point on $\overline{AB}$ with $2 \cdot AC = BC$. Let $D$ and $E$ be points on the circle such that $\overline{DC} \perp \overline{AB}$ and $\overline{DE}$ is a second diameter. What is the ratio of the area of $\triangle DCE$ to the area of $\triangle ABD$?


 
                    \begin{figure}[H]
                        \centering
                        \includegraphics[width=0.5\textwidth]{Problem/AMC_12_12A_2005_15_Problem_Figure_1.png}
                    \end{figure}
                    
$\textbf{(A) } \frac{1}{6} \qquad \textbf{(B) } \frac{1}{4} \qquad \textbf{(C) } \frac{1}{3} \qquad \textbf{(D) } \frac{1}{2} \qquad \textbf{(E) } \frac{2}{3}$


	\item (\textit{AMC 12 12A 2005}) Three circles of radius $s$ are drawn in the first quadrant of the $xy$-plane. The first circle is tangent to both axes, the second is tangent to the first circle and the $x$-axis, and the third is tangent to the first circle and the $y$-axis. A circle of radius $r > s$ is tangent to both axes and to the second and third circles. What is $r/s$?


 
                    \begin{figure}[H]
                        \centering
                        \includegraphics[width=0.5\textwidth]{Problem/AMC_12_12A_2005_16_Problem_Figure_1.png}
                    \end{figure}
                    
$\textbf{(A) } 5 \qquad \textbf{(B) } 6 \qquad \textbf{(C) } 8 \qquad \textbf{(D) } 9 \qquad \textbf{(E) } 10$


	\item (\textit{AMC 12 12A 2005}) A unit cube is cut twice to form three triangular prisms, two of which are congruent, as shown in Figure $1$. The cube is then cut in the same manner along the dashed lines shown in Figure $2$. This creates nine pieces. What is the volume of the piece that contains vertex $W$?


 
                    \begin{figure}[H]
                        \centering
                        \includegraphics[width=0.5\textwidth]{Problem/AMC_12_12A_2005_17_Problem_Figure_1.png}
                    \end{figure}
                    
$\textbf{(A) } \frac{1}{12} \qquad \textbf{(B) } \frac{1}{9} \qquad \textbf{(C) } \frac{1}{8} \qquad \textbf{(D) } \frac{1}{6} \qquad \textbf{(E) } \frac{1}{4}$


	\item (\textit{AMC 12 12A 2005}) Call a number "prime-looking" if it is composite but not divisible by $2$, $3$, or $5$. The three smallest prime-looking numbers are $49$, $77$, and $91$. There are $168$ prime numbers less than $1000$. How many prime-looking numbers are there less than $1000$?


 $\textbf{(A) } 100 \qquad \textbf{(B) } 102 \qquad \textbf{(C) } 104 \qquad \textbf{(D) } 106 \qquad \textbf{(E) } 108$


	\item (\textit{AMC 12 12A 2005}) A faulty car odometer proceeds from digit $3$ to digit $5$, always skipping the digit $4$, regardless of position. For example, after traveling one mile the odometer changed from $000039$ to $000050$. If the odometer now reads $002005$, how many miles has the car actually traveled?


 $\textbf{(A) } 1404 \qquad \textbf{(B) } 1462 \qquad \textbf{(C) } 1604 \qquad \textbf{(D) } 1605 \qquad \textbf{(E) } 1804$


	\item (\textit{AMC 12 12A 2005}) For each $x$ in $[0,1]$, define


 
 \[f(x) = \begin{cases}  2x, \qquad\qquad \mathrm{if} \quad 0 \leq x \leq \frac{1}{2}\\  2-2x, \qquad \mathrm{if} \quad \frac{1}{2} < x \leq 1.  \end{cases}\]


 Let $f^{[2]}(x) = f(f(x))$, and $f^{[n + 1]}(x) = f^{[n]}(f(x))$ for each integer $n \geq 2$. For how many values of $x$ in $[0,1]$ is $f^{[2005]}(x) = 1/2$?


 $\textbf{(A) } 0 \qquad \textbf{(B) } 2005 \qquad \textbf{(C) } 4010 \qquad \textbf{(D) } 2005^2 \qquad \textbf{(E) } 2^{2005}$


	\item (\textit{AMC 12 12A 2005}) How many ordered triples of integers $(a,b,c)$, with $a \geq 2$, $b \geq 1$, and $c \geq 0$, satisfy both $\log_{a}b = c^{2005}$ and $a + b + c = 2005$?


 $\textbf{(A) } 0 \qquad \textbf{(B) } 1 \qquad \textbf{(C) } 2 \qquad \textbf{(D) } 3 \qquad \textbf{(E) } 4$


	\item (\textit{AMC 12 12A 2005}) A rectangular box $P$ is inscribed in a sphere of radius $r$. The surface area of $P$ is $384$, and the sum of the lengths of its $12$ edges is $112$. What is $r$?


 $\textbf{(A) } 8 \qquad \textbf{(B) } 10 \qquad \textbf{(C) } 12 \qquad \textbf{(D) } 14 \qquad \textbf{(E) } 16$


	\item (\textit{AMC 12 12A 2005}) Two distinct numbers $a$ and $b$ are chosen randomly from the set $\{2, 2^2, 2^3, \ldots, 2^{25}\}$. What is the probability that $\log_{a}b$ is an integer?


 $\textbf{(A) } \frac{2}{25} \qquad \textbf{(B) } \frac{31}{300} \qquad \textbf{(C) } \frac{13}{100} \qquad \textbf{(D) } \frac{7}{50} \qquad \textbf{(E) } \frac{1}{2}$


	\item (\textit{AMC 12 12A 2005}) Let $P(x) = (x - 1)(x - 2)(x - 3)$. For how many polynomials $Q(x)$ does there exist a polynomial $R(x)$ of degree $3$ such that $P(Q(x)) = P(x) \cdot R(x)$?


 $\textbf{(A) } 19 \qquad \textbf{(B) } 22 \qquad \textbf{(C) } 24 \qquad \textbf{(D) } 27 \qquad \textbf{(E) } 32$


	\item (\textit{AMC 12 12A 2005}) Let $S$ be the set of all points with coordinates $(x,y,z)$, where $x$, $y$, and $z$ are each chosen from the set $\{0, 1, 2\}$. How many equilateral triangles have all their vertices in $S$?


 $\textbf{(A) } 72 \qquad \textbf{(B) } 76 \qquad \textbf{(C) } 80 \qquad \textbf{(D) } 84 \qquad \textbf{(E) } 88$


\end{enumerate}